% TOP_PROBLEM: {"id":30007518,"problem_number":"LOCAL-30007518","title":"Welfare costs of fluctuations","category_id":21,"category":{"id":21,"name":"economics","display_name":"Economics","description":"Open economic theory statements, published research agendas, and proposed scoped specifications; question provenance and historical priority are qualified per record.","slug":"economics","order_index":21,"created_at":"2026-10-08T00:00:00Z"},"status":"provenance_pending","external_url":null,"legacy_ids":[],"published":false,"statement":"Identify heterogeneous consumption-equivalent welfare costs of aggregate fluctuations under stated preferences and an explicit stabilization counterfactual.","economics":{"original_id":7,"field":"Business cycles and macro dynamics","kind":"identification","formalization_basis":"proposed_specification","source_status":"not_verified","priority_status":"not_established","priority_note":"No explicit primary open-problem statement matching this catalogue item was verified; supporting research is not evidence of first formulation.","earliest_verified_reference":null,"current_reference":{"authors":["Dimitris Georgarakos","Kwang Hwan Kim","Olivier Coibion","Myungkyu Shim","Myunghwan Andrew Lee","Yuriy Gorodnichenko","Geoff Kenny","Seowoo Han","Michael Weber"],"title":"How Costly Are Business Cycle Volatility and Inflation? A Vox Populi Approach","year":2025,"url":"https://www.iza.org/en/publications/dp/17675/how-costly-are-business-cycle-volatility-and-inflation-a-vox-populi-approach","doi":"10.3386/w33476","locator":"IZA Discussion Paper 17675 landing-page abstract, February 2025","evidence_summary":"The authors estimate households’ reported willingness to pay to eliminate fluctuations and find heterogeneous costs. The abstract reports findings; it does not state the proposed welfare-identification problem below."},"formalization_note":"New welfare-identification specification; the cited study provides estimates rather than an explicit open statement of this target."},"statusQualification":"Provenance pending: an explicit published statement of this exact open question was not verified. This is a catalogue-authored candidate specification, not a certified open conjecture. New welfare-identification specification; the cited study provides estimates rather than an explicit open statement of this target.","statusReviewedAt":"2026-10-08"}
% FIRST_POSTED: 2026-10-08
% ENTRY_KIND: statement_only
% !TeX program = lualatex
\documentclass[11pt]{article}
\usepackage{fontspec}
\usepackage[a4paper,margin=25mm]{geometry}
\usepackage{amsmath,amssymb,mathtools}
\usepackage{xurl}
\usepackage[hidelinks]{hyperref}
\newcommand{\sref}[2]{\hyperlink{source:#1:#2}{[S#2]}}
\setlength{\emergencystretch}{3em}
\begin{document}
% problemId:30007518
\section*{Welfare costs of fluctuations}\label{30007518}
\subsection{Definitions and mathematical statement}
Fix a population of household types \(i\), social weights \(\omega_i\ge0\) summing to one, discount factor \(\beta\in(0,1)\), and utility \(u_i(c,\ell)\) over consumption and labor. Let \((c_{it},\ell_{it})\) be allocations in the observed economy and \((\bar c_{it},\bar\ell_{it})\) those in a specified counterfactual that removes aggregate innovations while preserving idiosyncratic risk, taxes, and mean resources. Assume utility is increasing in consumption, discounted expectations are finite, and any compensation exists uniquely with \(\lambda_i>-1\). Define each type’s consumption-equivalent fluctuation cost \(\lambda_i\) by
\[
E\sum_{t\ge0}\beta^t u_i((1+\lambda_i)c_{it},\ell_{it})
=E\sum_{t\ge0}\beta^t u_i(\bar c_{it},\bar\ell_{it}).
\]
Identify or bound the distribution of \(\lambda_i\) and its socially weighted mean from household consumption, employment, beliefs, and elicited willingness to pay. Explicitly distinguish utility-based compensation from hypothetical survey payments: specify the assumptions under which reported payments identify \(\lambda_i\), and test framing, perceived feasibility, and income effects. The stabilization counterfactual must incorporate changes in unemployment risk and general-equilibrium prices rather than simply smoothing measured consumption ex post. Assess sensitivity to risk aversion, persistent uncertainty, incomplete insurance, and welfare weights. The exercise yields costs conditional on declared preferences and stabilization technology, not a preference-free universal welfare cost of cycles.

\par\textbf{Formulation and scope.} New welfare-identification specification; the cited study provides estimates rather than an explicit open statement of this target.

\par\textbf{Open-question provenance.} An explicit published open-question statement was not verified. Sources below are supporting literature and must not be treated as a first listing.

\par\textbf{Historical priority.} No explicit primary open-problem statement matching this catalogue item was verified; supporting research is not evidence of first formulation.

\subsection{Short English statement}
Identify heterogeneous consumption-equivalent welfare costs of aggregate fluctuations under stated preferences and an explicit stabilization counterfactual.

\subsection{Sources}
\begin{itemize}\raggedright
\item[\textnormal{[S1]}]\hypertarget{source:30007518:1}{} Dimitris Georgarakos, Kwang Hwan Kim, Olivier Coibion, Myungkyu Shim, Myunghwan Andrew Lee, Yuriy Gorodnichenko, Geoff Kenny, Seowoo Han, Michael Weber. 2025. How Costly Are Business Cycle Volatility and Inflation? A Vox Populi Approach. IZA Discussion Paper 17675 landing-page abstract, February 2025.
\url{https://www.iza.org/en/publications/dp/17675/how-costly-are-business-cycle-volatility-and-inflation-a-vox-populi-approach}.
DOI: 10.3386/w33476.
{\small\textit{Supporting or current literature; see evidence qualification. The authors estimate households’ reported willingness to pay to eliminate fluctuations and find heterogeneous costs. The abstract reports findings; it does not state the proposed welfare-identification problem below.}}
\end{itemize}
\end{document}
